\documentclass[12pt,a4paper]{article}
\usepackage[T2A]{fontenc}
\usepackage[utf8]{inputenc}
\usepackage[russian]{babel}
\usepackage{geometry}
\geometry{left=30mm,right=15mm,top=20mm,bottom=20mm}
\usepackage{setspace}
\onehalfspacing
\usepackage{indentfirst}
\setlength{\parindent}{1.25cm}
\usepackage{enumitem}
\setlist[enumerate]{leftmargin=*,itemsep=0.6em}
\begin{document}

\begin{center}
\textbf{\Large Анализ бахадуровской эффективности критериев согласия}
\end{center}

\section*{1. Цель работы}

Разработать численный метод оценки асимптотической эффективности критериев согласия по Бахадуру, реализовать его в виде нового типа эксперимента во фреймворке \texttt{pysatl-experiment}.

\section*{2. Задачи работы}

\noindent\textbf{1. Изучение теоретических основ эффективности критериев согласия по Бахадуру.}

\noindent 1.1. Изучить понятия точного и приближённого бахадуровских наклонов и подходы к их вычислению.

\noindent 1.2. Определить характеристики статистик, необходимые для численной оценки эффективности.

\noindent 1.3. Изучить применение информации Кульбака~--- Лейблера для оценки нормированной эффективности критериев согласия.

\medskip
\noindent\textbf{2. Проектирование численного метода и доработка фреймворка \texttt{pysatl-experiment}.}

\noindent 2.1. Разработать алгоритм численной оценки приближённого бахадуровского наклона с использованием ранее полученных предельных распределений статистик при нулевой гипотезе и выборок из альтернативных распределений.

\noindent 2.2. Изучить архитектуру фреймворка \texttt{pysatl-experiment} и определить необходимые изменения для реализации нового типа эксперимента.

\noindent 2.3. Спроектировать этапы вычисления, хранения и представления результатов оценки эффективности.

\medskip
\noindent\textbf{3. Программная реализация эксперимента \texttt{bahadur\_efficiency}.}

\noindent 3.1. Реализовать вычислительный модуль (воркер) для оценки бахадуровского наклона и эффективности критериев согласия.

\noindent 3.2. Реализовать настройку параметров эксперимента и интегрировать новый тип эксперимента в существующий фреймворк.

\noindent 3.3. Разработать тесты для проверки корректности работы вычислительного модуля.

\medskip
\noindent\textbf{4. Проведение вычислительных экспериментов и валидация результатов.}

\noindent 4.1. Провести численные эксперименты для выбранных критериев согласия и альтернативных распределений.

\noindent 4.2. Сравнить полученные численные оценки с известными аналитическими значениями.

\noindent 4.3. Исследовать точность, сходимость и вычислительные характеристики разработанного метода.

\medskip
\noindent\textbf{5. Анализ и оформление полученных результатов.}

\noindent 5.1. Проанализировать результаты вычислительных экспериментов и сформулировать выводы.

\noindent 5.2. Подготовить отчётную документацию, описание реализованного метода и результаты его проверки.

\end{document}
